\subsection{Infrastructure as a Service (IaaS) as the Experimental Substrate}
\label{subsec:iaas-substrate}

For the specific purpose of this thesis---instantiating and empirically validating the PANOPTES observability reference architecture---we adopt the Infrastructure as a Service (IaaS) model. This architectural decision is driven by the requirement for granular administrative access to the Kubernetes control plane, the Container Network Interface (CNI), Linux control groups (\texttt{cgroups}), and kernel-level instrumentation interfaces, which are partially or fully abstracted away in managed Platform as a Service (PaaS) Kubernetes offerings.

To guarantee scientific reproducibility and simulate a dedicated private IaaS environment free from multi-tenant cloud interference (the ``noisy neighbor'' effect) \cite{kuria_2024}, we employ a local multi-node infrastructure approach configured declaratively via \textbf{Ansible} and \textbf{Helm}. As detailed in Chapter~\ref{chap:experiment-implementation} (Section~\ref{subsec:virtualization-evolution} and \textbf{ADR-0006}), this substrate evolved across our iterative Design Cycle from VirtualBox/Multipass virtual machines to containerized multi-node Rancher K3s clusters (\textbf{K3d}) governed by strict per-node Linux \texttt{cgroup} hardware quotas ($2\,\text{vCPUs}$ and $4096\,\text{MiB}$ RAM per node), providing a deterministic testbed for evaluating cloud-native observability under controlled fault injection.
